% Generated draft: requires mathematical and source review.
\documentclass[11pt]{article}
\usepackage[margin=25mm]{geometry}
\usepackage{fontspec}
\setmainfont{Latin Modern Roman}
\usepackage{amsmath,amssymb,mathtools,unicode-math}
\setmathfont{Latin Modern Math}
\usepackage{enumitem,needspace,xurl,hyperref}
\setlength{\parindent}{0pt}
\setlength{\parskip}{5pt}
\setlength{\emergencystretch}{3em}
\newcommand{\N}{\mathbb{N}}
\newcommand{\Z}{\mathbb{Z}}
\newcommand{\Q}{\mathbb{Q}}
\newcommand{\R}{\mathbb{R}}
\newcommand{\C}{\mathbb{C}}
\newcommand{\F}{\mathbb{F}}
\newcommand{\A}{\mathbb{A}}
\newcommand{\PP}{\mathbb{P}}
\newcommand{\E}{\mathbb{E}}
\newcommand{\eps}{\varepsilon}
\newcommand{\dd}{\,\mathrm{d}}
\begin{document}
\section*{UnsolvedMath ID: 30005243}
\textbf{Source problem code:} catalogue entry
\subsection{Definitions and mathematical statement}
\textbf{Source-limited formulation.} The supplied record concerns high-frequency wave propagation near an obstacle. An incident ray is called tangent when its direction at contact lies in the tangent space of the obstacle boundary. In the supplied description, a creeping ray travels along the boundary following a boundary geodesic and radiates surface-diffracted rays into the shadow region. The shadow and transition regions are mentioned but are not mathematically defined in the record.

The source describes part of a circular front of directly propagated singularities and a curved front of singularities reflected from the obstacle according to Snell's law. It further describes an incident tangent ray as splitting into an undeviated branch along the shadow boundary and a branch that travels along the obstacle boundary, which acts as a secondary source by radiating into the shadow region. The source asserts exponential decay of surface intensity with distance along the ray and describes creeping rays as reaching the transition and shadow regions. These are supplied descriptive claims; their precise hypotheses, scales, and quantitative meanings are absent.

\textbf{Unrecoverable mathematical target.} The source ends with the incomplete assertion: ``A result stated by Keller (and proved by [Harg\'e \& Lebeau, 1994] in case of $C^\infty$ boundaries) states that rays can carry a non-negligeble amount of energy when''. The condition following ``when'' is missing. Consequently, neither this asserted result nor a quantified question about diffractive propagation can be recovered from the supplied text. In particular, the record supplies no governing wave equation, ambient dimension, obstacle geometry or convexity assumptions, boundary condition, incident data, high-frequency parameter, definition of energy or intensity, meaning of non-negligibility, or target estimate.

The registry marks the statement as \texttt{unrecoverable}. Its supplied statement-verification note says that the report treats this heading through sharp results and counterexamples rather than stating an unresolved mathematical problem. This note is supplied curation, not an independently verified conclusion. No precise conjecture or theorem is asserted in this draft.
\subsection{Short English statement}
The record describes tangent rays, creeping propagation, and radiation into an obstacle's shadow, but its final energy assertion is cut off before the essential condition. A precise mathematical problem cannot be reconstructed. The supplied \texttt{partially\_solved} label is a curation status, and the record itself marks the statement as \texttt{unrecoverable}.
\subsection{Sources}
\begin{itemize}
\item Local UnsolvedMath registry: \texttt{data/unsolvedmath/problems.json}; supplementary registry: \texttt{catalog/problems.json}. Supplied permanent ID: \texttt{30005243}; exact source problem code: \texttt{catalogue entry}; title: \emph{Diffractive High-Frequency Wave Propagation}. The record belongs to the Oberwolfach Reports Open Problems collection, records proposed year 2022, and names no proposer.
\item Supplied report citation: \emph{At the Interface between Semiclassical Analysis and Numerical Analysis of Wave Scattering Problems} (2023), Oberwolfach Reports, pp.\ 2511--2587, DOI \texttt{10.4171/owr/2022/43}. Source URL and supplied literature-source link, labelled ``Oberwolfach Report 43/2022'': \url{https://doi.org/10.4171/owr/2022/43}. The supplied background identifies Oberwolfach Reports as records of mathematical workshops and their open-problem discussions.
\item Supplied literature-source entry: ``Ivanovici--Lebeau, Comptes Rendus Math\'ematique'', \url{https://doi.org/10.1016/j.crma.2017.05.011}. No article title or precise result location is supplied.
\item Attribution within the truncated source: a result stated by Keller and reportedly proved by Harg\'e \& Lebeau (1994) for $C^\infty$ boundaries. No complete bibliographic reference is supplied for this attribution, and the asserted result's concluding condition is missing.
\item Literature assessment dated exactly \texttt{2026-08-22}: supplied dated curation, not independently verified. The supplied status is \texttt{partially\_solved}. The assessment says that creeping rays, surface radiation, and Poisson-spot effects have rigorous descriptions in important smooth convex cases, while the report emphasizes losses and incomplete quantitative control for general boundaries. The background describes this as dated literature triage derived from supplied OWR curation and requests independent verification of time-sensitive status and cited claims. No such verification is claimed here.
\end{itemize}
\end{document}
